\begin{pastexam}[label={px:ntu-108-ee-dsb-01}]{臺大電機所 108 學年度（B）第 1 題}
（是非；答對 4 分，答錯倒扣 4 分）

A binary heap may not be a complete binary tree.
\end{pastexam}

\begin{pastsoln}
答案是 (B)，錯。

二元堆積的定義有兩個條件：形狀是完全二元樹，而且每個節點的鍵不小於（最小堆積是不大於）它的
子節點（資結書第~8.1~節）。 形狀這個條件是定義的一部分，不是碰巧成立的性質：正因為是完全
二元樹，才能把節點依層序放進陣列 $A[1..n]$，用 $\lfloor i/2 \rfloor$、$2i$、$2i+1$ 找父節點和
子節點，中間不留空格；也因為是完全二元樹，高度才是 $\lfloor \lg n \rfloor$，插入和取出最大值
才是 $O(\lg n)$。 一棵二元樹如果不是完全二元樹，就算滿足堆積的大小關係，也不叫二元堆積。

題目說「可能不是完全二元樹」，和定義矛盾，所以錯。
\end{pastsoln}
